\documentclass[12pt]{article}
\usepackage[utf8]{inputenc}
\usepackage[LGR, T1]{fontenc}
\usepackage[greek.ancient,english]{babel}
\usepackage{alphabeta}

\usepackage[letterpaper, margin = 1in]{geometry}
\usepackage{fancyhdr}
\usepackage{titling}
\usepackage{titlesec}
\usepackage{hyperref}
\usepackage{lastpage}
\usepackage{amssymb}
\usepackage{setspace}

\makeatletter
\newcommand{\hangpara}[2]{\hangindent#1\hangafter#2\noindent}
\newenvironment{hangparas}[2]{\setlength{\parindent}{\z@}%
  \everypar={\hangpara{#1}{#2}}}{\par}
\makeatother

\providecommand{\medium}{\normalsize}

\title{Boredom in Adorno, Baudelaire, and Benjamin}
\author{Ian Thomas White}
\date{}

\makeatletter
\renewcommand{\maketitle}
{
\begin{flushleft}
\textbf{{\large\@title}}
\newline
\@author
\end{flushleft}
}
\makeatother

\titleformat*{\section}{\large\bfseries}
\titleformat*{\subsection}{\medium\bfseries}
\titleformat*{\subsubsection}{\small\bfseries}

\pagestyle{fancy}
\fancyhf{}
\setlength{\headheight}{14.5pt}
\lhead{Ian Thomas White}
\chead{Boredom in Adorno, Baudelaire, and Benjamin}
\rhead{\thepage \ of \pageref{LastPage}}

\doublespacing
\widowpenalty=100000
\clubpenalty=100000

\begin{document}

\thispagestyle{empty}

\begin{singlespacing}
\maketitle
\tableofcontents
\end{singlespacing}

\newpage

\section{Introduction}
The consistently strained, perhaps dire, tone of Adorno's writing is hard to miss, and Adorno is relatively explicit, at times, about the intention to shock that such a tone is meant to carry out. This emphasis on shock links Adorno to Baudelaire, especially via Benjamin's interpretation of the poet. Adorno's dependence on Baudelaire and Benjamin is especially clear in Adorno's late opus \textit{Aesthetic Theory}, insofar as Adorno there attempts to deal with modernity and modern art under the two writers' influence. Despite his understanding of them, however, Adorno in his aesthetics fails to take into account the crucial role that \textit{ennui}, boredom, takes in Baudelaire's and Benjamin's thought. By rooting his aesthetics in terms derived from Baudelaire while leaving out the concept of boredom, Adorno ends up missing what would otherwise seem to be, more generally, an important aspect of modernity and what could provide a useful structure for his own theory. The three major aims of this paper, then, are as follows: 1) to trace the importance of the Baudelaire-Benjamin influence on \textit{Aesthetic Theory}; 2) to show, unlike those two, Adorno's failure to account for \textit{ennui}, boredom; 3) to describe the relevance that failure has within and for Adorno's theory.

\section{Adorno's Style}
Adorno's difficulty-bearing bombast is both known and apparent. Briefly, the translator of \textit{Aesthetic Theory} says that the work is ``unparalleled'' with regard to its paratactic nature and the ``density'' that nature entails, that it ``necessarily tends toward the apodictic'', and that it is ``inimical to exposition''.\footnote{Robert Hullot-Kentor, introduction to \textit{Aesthetic Theory}, by Theodor Adorno, trans. and ed. by Robert Hullot-Kentor (Minneapolis, MN: University of Minnesota Press, 1997), xvi.} Similarly, Sianne Ngai says \textit{Aesthetic Theory} has an ``unusually dense and protracted form''.\footnote{Sianne Ngai, \textit{Our Aesthetic Categories: Zany, Cute, Interesting} (Cambridge, MA: Harvard University Press, 2012), 101.} Adorno is, in accordance with his style, not shy about his style. He happily, habitually takes on the essay form because it ``arouses resistance'',\footnote{Theodor Adorno, ``The Essay as Form'' in \textit{Notes to Literature, Volume 1}, ed. by Ralph Tiedemann, trans. by Shierry Weber Nicholsen (New York: Columbia University Press, 1991), 3.} and he approves of the use of foreign words due to their ``dangerous'' quality.\footnote{Theodor Adorno, ``On the Use of Foreign Words'' in \textit{Notes to Literature, Volume 2}, ed. by Ralph Tiedemann, trans. by Shierry Weber Nicholsen (New York: Columbia University Press, 1992), 291.} Whether or not one takes seriously Adorno's caveat that ``dialectics does not give any instructions for the treatment of art'',\footnote{Theodor Adorno, \textit{Aesthetic Theory}, trans. and ed. by Robert Hullot-Kentor (Minneapolis, MN: University of Minnesota Press, 1997), 140.} \textit{Aesthetic Theory} provides a wall of text built of predicate-nominative-based bricks, the whole of which might be said to be aimed at skulls.  

Of course, Adorno isn't being difficult for the sake of being difficult; i.e., even if he is invested in ``violations of the orthodoxy of thought'',\footnote{Adorno, ``Essay as Form'', 23.} he's not just trying to provide an occasion for dismounting a reader's apathy by providing a simple contrast to the easier flows of other texts. Rather, Adorno's style programmatically reflects conceptions of modernity from Baudelaire, especially via Benjamin,\footnote{The translator of \textit{Aesthetic Theory} also claims that it is ``as a whole, a response'' to Benjamin's ``Artwork in the Age of Mechanical Reproduction'', Hullot-Kentor,  introduction, xvi.} conceptions which center on the concept of the fragment. Adorno's extollation of the essay form, that is, seems dependent on his reading of Baudelaire. Of his own use of the essay, Adorno claims,   

\begin{quote}
\begin{singlespace}
Its [the essay's] self-relativization is inherent in its form: it has to be constructed as though it could always break off at any point. It thinks in fragments, just as reality is fragmentary, and finds its unity in and through the breaks \dots  An unequivocal logical order deceives us about the antagonistic nature of what that order is imposed upon.\footnote{``Essay'', 16.}
\end{singlespace}
\end{quote}

\noindent In \textit{Aesthetic Theory}, furthermore, Adorno claims something similar of Baudelaire's writing,

\begin{quote}
\begin{singlespace}
Baudelaire, the apologist of form no less than the poet of the \textit{vie moderne}, expressed this [impossibility of homeostasis between the universal and the particular] in the dedication of \textit{Le spleen de Paris} when he wrote that he can break off where he pleases, and so may the reader, `for I have not strung his wayward will to the endless thread of some unnecessary plot'.\footnote{\textit{Aesthetic Theory}, 224. The editors cite Charles Baudelaire, ``To Arsene Houssaye'', in \textit{The Poems in Prose}, trans. Scarfe (London, 1989), 25.}
\end{singlespace}
\end{quote}

\noindent In both passages Adorno highlights: 1) the centrality of fungible fragments for his and Baudelaire's project, 2) the motivating force that ``form'' takes in organizing those fragments, and 3) the danger of misleading the reader in some way when fragments aren't utilized. Most importantly, the textual fragments have a reflective-like relationship to some outside-text reality, the \textit{vie moderne}. As will be seen, given the paratactic way each section of \textit{Aesthetic Theory} is stitched together,\footnote{For comparison, see also Plato's rejection of a very similar paratactic ``stitching'' method for use in rhetoric/poetry in the \textit{Phaedrus}. For one translation: Plato, \textit{Phaedrus}, trans. by Alexander Nehamas and Paul Woodruff (Indianapolis, IN: Hackett Publishing, 1995), 61-62.} it seems fair to say Adorno is trying to transport the essay form into \textit{Aesthetic Theory}. The book, then, too, trades on a form born of Baudelaire.   

Benjamin's interpretation of Baudelaire emphasizes the same set of concerns. He focuses on the fragment and its relation to reality insofar as ``Baudelaire holds in his hands the scattered fragments of genuine historical experience''.\footnote{Walter Benjamin, ``Some Motifs in Baudelaire'' in \textit{Illuminations: Essays and Reflections}, trans. by Harry Zohn, and ed. by  Hannah Arendt (New York: Schocken Books, 1968), 185.} Benjamin also relates fragments to a totalized form, since ``multiple fragments [can be] \dots  assembled under a new law''.\footnote{Ibid., 234.} Finally, he says that Baudelaire's method comes undistorted, the reverse, at least, of Adorno's concern with misleading the reader: ``the experience \dots  presented itself to Baudelaire's eyes in its undistorted version in the figure of his reader''.\footnote{Ibid., 157.} As with his style, Adorno isn't necessarily coy in \textit{Aesthetic Theory} about the fragment's significance to modern life, such that he often mentions its role in aesthetics and society.\footnote{E.g., Adorno, \textit{AT}, 90, ``Its [i.e., art's] highest products are condemned to a fragmentariness that is their confession that even they do not possess what is claimed by the immanence of their form.''} My point, here, is to show Adorno's explicit inheritance of the fragment from Baudelaire and Benjamin, in order to (later) put in relief what Adorno leaves out from this inheritance.  

The fragment, as constitutive of modern life is, for Adorno,\footnote{Adorno at least implicitly includes the ``reality'' during which he was writing within modernity, e.g., Adorno, \textit{AT}, 20, ``\dots  incomparable power in Edgar Allan Poe, truly a beacon for Baudelaire and all modernity''.} linked to the concept of ``shock''. The fragment has a penchant for ``shattering''.\footnote{Adorno, ``Essay as Form'', 20.} Moreover, aesthetics is rooted in the fragmentary form, since ``art that makes the highest claim compels itself beyond form as totality and into the fragmentary''.\footnote{Adorno, \textit{AT}, 147.} Adorno links the concept of form more generally to shock value, as well: ``Technical critique of highly technical works obviously disregards a conception of composition that asserts the principle of shock as fundamental and actually transfers the unity of the composition into the irrational suspension of what traditional style called logicality, unity.''\footnote{Ibid., 215.} Finally, Adorno points to the relationship between shock and a potential viewer/reader. He claims, ``the most extreme shocks and gestures of alienation of contemporary art \dots  are nearer than they appear to be by virtue of historical reification. What is considered to be intelligible [i.e., the non-shocking] to all is what has become unintelligible''.\footnote{Ibid., 183.} For Adorno, then, the fragment and shock are, through the concept of ``form'', linked concepts, perhaps to a certain degree causally. They are constitutive of modern life and modern art insofar as they comprise an aggregate form and aren't de facto misleading.\footnote{E.g., ibid., 45, ``The fragment is that part of the totality of the work that opposes totality.''}  

Adorno inherits the fragment-shock link from Benjamin and Baudelaire. As is well known, Benjamin locates much of his conception of modernity in the shock. For Benjamin, ``the question suggests itself how lyric poetry [of which Baudelaire is paramount] can have as its basis an experience for which the shock experience has become the norm''.\footnote{Benjamin, ``Motifs'', 162.} This particular inheritance, though, opens up a new link to Adorno's text. Namely, Benjamin's concept of shock is directly related to the repetitive nature of modern life because ``the hand movement of the work at the machine \dots  repeats that [preceding] gesture exactly'', and there is a ``jolt in the movement of the machine''.\footnote{Ibid., 177.} The factory machine, an archetypal image of modernity, entails a series of repetitive actions, each of which sends a jolt. Baudelaire's classic poem of a man struck by the figurative lightning of the sight of a mourning woman makes shock no less central to his aesthetics.\footnote{Charles Baudelaire, ``À une passante'' in \textit{Fleurs du Mal (1861 Edition)}, last modified 2018, \url{https://fleursdumal.org/poem/224}. Benjamin, ``Motifs'', 168-169, analyzes this poem, as well.}  

\textit{Aesthetic Theory}, too, as a text, is obviously repetitive. Readers might, in fact, ``rankle at'' the ``repetitiveness''.\footnote{Hullot-Kentor, introduction, xvii.} Adorno repeats things within and between his sections, and the implication is that, if the essay is a fragment, \textit{Aesthetic Theory} consists of a series of shock-inducing fragments, directly relating it to the modern predicament. Thus, Adorno's reliance on Baudelaire and Benjamin for his own stylistic method (the ``form'') and for an important subset of his conceptual apparatuses (the ``content'') provides refined insight into the textual wall \textit{Aesthetic Theory} presents. Rather than mere difficulty-for-difficulty, the bombast of \textit{Aesthetic Theory} combines Benjamin's conception of repetition with the aesthetic form of fragmentation,\footnote{Though Benjamin, ``Motifs'', 176, does call the laborer's experiences ``isolated''.} with the aim to shock the reader, rather than to mislead her with some classically-timed unity. This is further verified when considering Adorno's project and goal for aesthetic theory more generally, i.e., that ``aesthetics must no more lag behind art than behind philosophy''.\footnote{Adorno, \textit{AT}, 344.} Adorno, in order to best reflect the reification of modern life has his conceptual apparatus reflect that life's repetitive fragmentation.\footnote{This doesn't mean \textit{Aesthetic Theory} is, for Adorno, itself an ``artwork''. The work rather reflects ``the philosophical insight that fact and concept are not polar opposites but mediated reciprocally in one another'', ibid., 343.} Adorno's paratactic style is meant to keep the reader shuddering, and each repetition, never an exact replica, seems to want to keep the reader on her toes, demanding constant attention. Adorno's combination of these effects is in direct lineage of Benjamin and Baudelaire.\footnote{To be clear, I'm not trying to say that Adorno is merely being derivative. Adorno, of course, in addition to his fragment-repetition combination has intricate developments from Benjamin and Baudelaire in terms of what actually constitutes shock, e.g., the relationship between ``shudder'', form, and beauty (ibid., 20, 79).}

\section{Boredom in Baudelaire and Benjamin}
I will briefly make an interlude to focus directly on Benjamin and Baudelaire, namely on their presentation of a set of phenomena I'll label, for convenience, boredom. Benjamin's description of shocking repetition, the series of unrelated hand movements, in both the factory and the gambling hall hinges on the label of ``drudgery'',\footnote{Benjamin, ``Motifs'', 177.} and he categorizes an inveterate gambler's testimony to the non-``boring'' nature of the game as ``phantasmagoria''. Benjamin also emphasizes Baudelaire's insistence on boredom. He quotes in full Baudelaire's prose poem ``A Lost Halo'', which solidifies the link between lyric ``dignity'' and ``boredom'' (Baudelaire's terms),\footnote{Ibid., 192-193.} and he earlier, explicitly notes Baudelaire's link between the ``dull'' and ``beauty'' (Benjamin's terms).\footnote{Ibid., 190.}

Baudelaire himself clearly locates his sense of modern life and of aesthetics in the experience of boredom. What amounts to a preface for \textit{Fleurs du Mal}, ``Au Lecteur'' is a poem which features several lists of near synonyms (of various flaws, sins, adjectives, etc.) and the poem itself speeds through a serialized list of semi-redundant, hyperbolic, \textit{maudit} metaphors.\footnote{Charles Baudelaire, ``Au Lecteur'' in \textit{Fleurs du Mal (1861 Edition)}, last modified 2018, \url{https://fleursdumal.org/poem/099}. I've included the whole poem at the end as an appendix} For the climax of the poem, Baudelaire inserts an entire stanza listing, in a rigid meter,\footnote{As is most of the poem, but the list form highlights it.} wild animals and their cries. The conclusion of this wild ``ménagerie'' is that ``II en est un plus laid, plus méchant, plus immonde / \dots  Et dans un bâillement avalerait le monde // C'est l'Ennui!''. \textit{Ennui}'s world-engulfing yawn implies the concept is forcibly prevalent to the modern experience, while that yawn's conditional, i.e. unrealized, and tired, nature implies the continuing dissatisfaction that attends the affect. In other words, ``Au Lecteur'' operates on the tension between ``drudgery'' and the shocking as a function of (metered) repetition---overlong lists of beasts and ever unsatisfied, earth-sized boredom.  

Crucially, moreover, \textit{Ennui} in the poem functions to refine the nature of subjecthood and aesthetics in the poem. Before the final stanza  a vague ``nous'' operates as the poem's primary, lyric subject, as if referring to ``humanity''. Upon \textit{Ennui}'s arrival, though, this ``nous'' explodes into a ``tu'' separated from the narrator (an implied yet hidden ``je''). The explosion has two effects: 1) it forces a reinterpretation of the former ``nous'' as a set of actual, experiencing subjects, 2) it subsumes newly born subjects to the poem's understanding of the modern life and its equivocal \textit{Ennui} (so much that the ``je'' never fully shows up). Furthermore, by calling the new reader and itself a ``hypocrite'', the narrator indicates the bourgeois reader's desire for something lyrical and beautiful (even somewhat satisfying it in an underhanded way),\footnote{As Benjamin, ``Motifs'', 156 says, Baudelaire ``has become a representative of a [lyric] genre''.} as well as the fact that the lyric form and its consequent desire has become unbearably dull. The shock of the narrator's name-calling/self-indictment and wild beasts, with the \textit{Ennui} climax, indicates the narrator takes the \textit{Ennui}\-shock relationship to be fundamental to aesthetic experience itself, all of which is expressed through, and curtailed by, a constant, precise lyric meter.

\section{Boredom in Adorno}
Despite what appears to be the centrality of boredom for Baudelaire and Benjamin and despite Adorno's inheritance from those two writers, boredom features relatively little in \textit{Aesthetic Theory}. On the contrary, Adorno seems almost to actively excise boredom from both aesthetic objects and aesthetic experience. Boredom explicitly shows up only twice in the whole work. It first appears in the ``Coherence and Meaning'' section during a discussion on the antagonism between conceiving the artwork as, on the one hand, ``authentic'' (i.e., a \textit{l'art pour l'art} aesthetic conception) and as, on the other hand, ``relevant'' (i.e., as subordinate to extra-aesthetic cultural values). Adorno claims that ``what ranks as problematic [for some others] in the aesthetic sphere has its origin here [in the relative loss of the artwork's cultural import]; the remainder became the plunder of boredom''.\footnote{Adorno, \textit{AT}, 158.} What's important to note here is that ``boredom'' specifically seems to find its place outside of, and detracts from, the ``aesthetic sphere'', or at least that sphere's leftovers  

Adorno even mentions in this context, i.e. in the same paragraph, ``art's paratactic logicality'' with regard to the relationship between the artwork's elements and the whole of the artwork.\footnote{Ibid, 157. The whole being a ``negative harmony'' that stands in ``dissonant relation'' to its elements.} Given his links to Baudelaire and Benjamin, that Adorno mentions both a paratactic form\footnote{This is one of only a few times Adorno himself uses ``paratactic''.} and boredom in the same breath seems more than coincidental; however, if Adorno were strictly following them here, the mention of art's paratactic form would entail the boredom. The two are not explicitly linked here, though. Whereas Baudelaire both maintains and ironically undercuts the overarching yawn of \textit{Ennui}, Adorno brushes off the image of an ``entire wholeness'' from ``traditional aesthetics'' as a mere ``mistake''.\footnote{Ibid. This isn't to say that Adorno unfairly privileges the elements over the whole. Adorno, of course, conceives of the whole-elements relationship dialectically, i.e., as mutually constitutive, as the quote here displays.} Even though the unity of an artwork may be dialectically significant for Adorno, he implies that a world-scale unity, unlike Baudelaire's \textit{Ennui}, isn't a positive or negative motivating concept for aesthetics.\footnote{The point is that Adorno doesn't take this seriously for his own, present aesthetic theory in the same way that Baudelaire does, even if the mistakes of former aesthetics are still historically/genealogically significant. Nor does this ``mistake'' seem negatively motivating here.} Furthermore, artworks that fail to adequately succeed in finding ``equilibrium'' of these paratactic elements, artworks that fail to ``gel'', just ``impotently dissipate''\footnote{Ibid.}---Adorno rejects boredom as motivating for either successful or failed artworks.   

The other instance of boredom in \textit{Aesthetic Theory} occurs almost at the end of the book, in one of its rare moments of reverie:  

\begin{quote}
\begin{singlespace}
Music, whether it is played in a cafe or, as is often the case in America, piped into restaurants, can be transformed into something completely different, of which the hum of conversation and the rattle of dishes and whatever else becomes a part. \dots  A medley is sometimes made up of parts of artworks, but through this montage the parts are fundamentally transformed. Functions such as warming people up and drowning out silence recasts music as something defined as mood, the commodified negation of the boredom produced by the grey-on-grey commodity world. The sphere of entertainment, which has long been integrated into production, amounts to the domination of this element of art over all the rest of its phenomena. \dots  Yet someone sitting in a cafe who is suddenly struck by the music and listens intensely may feel odd to himself and seem foolish to others. In this antagonism the fundamental relation of art and society appears.\footnote{Ibid., 253.}
\end{singlespace}
\end{quote}


Boredom is here again paired not (directly) with the aesthetic realm, but with the real, ``commodity world'', i.e. the realm of reified and non-autonomous objects. Of course, for Adorno, ``negation'' of boredom would not likely mean the complete exclusion of boredom, but the negation here appears to be limited in nature.\footnote{For another example of ``negation'', ibid., 15: ``The taboo on the sensual ultimately encroaches on the opposite of pleasure because, even as the remotest echo, pleasure is sensed in its specific negation.'' Here pleasure and its opposite are certainly linked and understanding one requires the other, but this is not the same as saying they constitute or determine each other in a substantive way.} In this scene, boredom \textit{per se} is not constitutive of the experience, aesthetic or otherwise, of the listener so described, despite boredom's constitutive nature according to Baudelaire and Benjamin.  

Three other important distinctions pertain to the boredom described by Adorno and that of Baudelaire. First, while Adorno, like Baudelaire, maintains the engulfing danger of the boredom-based phenomenon (that philistine pleasure), unlike Baudelaire, ``this element'' for Adorno actually succeeds in ``dominating'' the other elements.\footnote{I take the genitive in ``the domination of this element'' as (grammatically) subjective rather than objective.} While Baudelaire's \textit{Ennui} is stuck in an overarching yet unrealized conditional, Adorno's reaction to boredom actually succeeds in overriding other concerns (to negative effect, of course). Even if the pleasure from the medley is somewhere along the line rooted in boredom, the pleasure-based element would override even the boredom itself. Second, the nature of the boredom is different here than in Baudelaire. The paratactic-like form, the medley/montage, is precisely the phenomenon here that overrides the boredom of modern reality, rather than constitute it; contrarily, the ``grey-on-grey'' form of actual boredom implies a constant synchronic sameness. This synchronic form is distinctly opposed to the diachronic repetitive boredom in Benjamin and Baudelaire.  

Third, and most importantly, boredom and its related pleasure is completely cordoned off from a more legitimate aesthetic experience. The listener's up-pricked ears do stem from a shock---the listener is ``suddenly struck''. This shock, though, is not funded by the paratactic form provided in this context. Furthermore, rather than recouping the underlying boredom of the lounge music, per its commodified nature, the epiphany results in the listener's newfound ``intensity''; the listener feels ``odd'' (the opposite of dull) and becomes marked, a spectacle (the opposite of background). The ``fundamental antagonism'' here between the artwork as cultural background and as the listener's solitary epiphany (a reprise of the antagonism in the first boredom quote between \textit{l'art pour l'art} and cultural value) places boredom squarely on one side of the debate: the non-aesthetic, reified world.  

The closest Adorno gets, I think, to Benjamin's link between repetition and boredom is in the ``Paralipomena'' of \textit{Aesthetic Theory}. After bringing up Benjamin's link to ``industrial procedures'', Adorno claims that `` the repetition of identical rhythms and the repetitive manufacture of an identical object based on a pattern, at the same time contain a principle antithetical to the new.''\footnote{Ibid., 272.} It may be possible to read boredom into ``the antithetical to the new'', but Adorno is at best being coy about the experience that industrial techniques cause.

\section{The Potential Importance of Boredom for \textit{Aesthetic Theory}}
On a relatively trivial level, Adorno's oversight or rejection with respect to boredom creates a simple consistency problem for \textit{Aesthetic Theory}. If Adorno's goal is indeed for aesthetics not to ``lag behind art''\footnote{Ibid., 344.} and to ``confront itself with its object'',\footnote{Ibid., 333.} Baudelaire's take on \textit{Ennui}, which would ostensibly be at least a subset of the objects that aesthetics is concerned with (and from a fundamental witness to modernity and modern aesthetics\footnote{In Adorno's own estimation. Adorno calls Baudelaire a ``herald of modernism'', ibid., 133 and says Baudelaire was the first to articulate a theory of the modern, ibid., 21.}), should be given more regard in \textit{Aesthetic Theory}. Nor, though it's outside the scope of this paper to demonstrate, are the aesthetic objects trading in repetitive boredom and shock confined to the poems of Baudelaire.  

More substantively, though, I think that the concept of boredom as manifest in Baudelaire can help solve an aporia in Adorno's aesthetics that centers on the nature of subjecthood. With regard to subjecthood as it manifests in an artwork, Adorno claims the following:  

\begin{quote}
\begin{singlespace}
Art's linguistic quality gives rise to reflection over what speaks in art; this is its veritable subject, not the individual who makes it or the one who receives it. This is masked by the lyrical "I" \dots  But this subjectivity is on no account identical with the I that speaks in the poem. \dots  The part played by the empirical I is not, as the \textit{topos} of sincerity would have it, the locus of authenticity \dots  Only as things do artworks become the antithesis of the reified monstrosity. Correspondingly, and this is key to art, even out of so-called individual works it is a We that speaks and not an I---indeed all the more so the less the artwork adapts externally to a We and its idiom.
\end{singlespace}
\end{quote}

\noindent Adorno here conceives\footnote{The following paragraph is a paraphrase of my ``Response Paper 4'' for this class.} of three distinct, potential subjects that an artwork might manifest: the creator-level, the narrator-level, and the spirit-level. The first level he dismisses as trivial, and the second is falsely over-empirical (and linked to \textit{l'art pour l'art} conceptions of the artwork). The primary operant subject in an artwork, then, pertains to the level of spirit. This subject is a collective ``we'' because it is an actual, created object that partakes from the same society that creates it---it expresses, to use Adorno's term, the qualities of the world it comes from. Furthermore, this subject, as spirit, is transcendental insofar as it is abstracted from individuals, and it thereby participates in determining social relations as modeled on the system of exchange operant in society.\footnote{See Theodor Adorno, ``On Subject and Object'' in \textit{Critical Models}, trans. by Henry Pickford (New York: Columbia University Press, 1999), 248.}  

The problem with respect to boredom, though, starts with what Adorno rejects initially as trivial: the creator-level subject. For this first-level subject, Adorno ascribes the quality opposite to that which he ascribes the transcendental subject; he ascribes to it necessary singularity rather than necessary plurality:  

\begin{quote}
\begin{singlespace}
Collective modes of production by small groups are already conceivable, and in some media even requisite; monads are the locus of experience in all existing societies. Because individuation, along with the suffering that it involves, is a social law, society can only be experienced individually. The substruction of an immediately collective subject would be duplicitous and would condemn the artwork to untruth because it would withdraw the single possibility of experience that is open to it today.\footnote{Adorno, \textit{AT}, 259.}
\end{singlespace}
\end{quote}

\noindent The question can be stated, then: How does a necessarily plural subject derive from a necessarily singular subject? This question seems to be one of the principle aporiae of \textit{Aesthetic Theory} and is noted, in various guises, as such and as still not fully solved.\footnote{See Ngai, \textit{Categories}, 105-109 for the at least related problem of sociality-via-antisociality. She argues for its unsolved nature and provides a potential solution different than mine.}  

Baudelaire's ``Au Lecteur'' provides a potential solution. As mentioned previously, before the poem's stanza explosion of ``mon semblable'' from the narrator, a ``nous'' is the primary lyrical subject of the poem, and this subject initially goes through a series of contorted metaphors about the disgusting qualities of the modern individual. Like Adorno's warning about a ``collective subject'', this ``nous'' is superficial, ``duplicitous'', insofar as it doesn't provide criteria for what constitutes the ``nous'' and insofar as the metaphors are self-consciously (and, of course, hypocritically) hyperbolic. With the onslaught of \textit{Ennui} and thereby of the ``mon semblable'', with the ``tu'' and hidden ``je'', the poem claims that \textit{Ennui} is in fact experienced individually, as Adorno says is true of society. At the same time the poem itself subsumes both subjects under \textit{Ennui} as a collective, real world (engulfing) problem. In other words, ``Au Lecteur'' provides a structure for rejecting a spurious ``we'' and then for moving between individuated subjects and a collective, reified world---precisely the structure that \textit{Aesthetic Theory} needs in order to solve its aporia. Benjamin links the poem and the factory, the lyrical I and the factory worker, in a no less felicitous way. Modern boredom, as found in the factory, the gambling hall, and the metered poem, seems to be a solution, if a bleak one, to the disjuncture between individual and collective.

\section{Conclusion}
\textit{Aesthetic Theory} is a difficult and repetitive text. The repetition inherent to it, though, never relents its requirement for concentration---it is never boring. Under some (even non-trivial) criteria, constant attention might be beneficial, but \textit{Aesthetic Theory}'s constant need to keep the reader on edge, despite keying into a crucial aspect of modernity, shock, seems to miss another crucial aspect, boredom. The structure of boredom, via Benjamin and Baudelaire, in addition to being constitutive of modernity, as an individuated experience that threatens yet fails to engulf, can provide a structure for linking individual and collective subjects. This structure would then help solve what seems to be an outstanding aporia for \textit{Aesthetic Theory}, and it wouldn't, however, entail a theoretical refocus on some subjective response to art, which Adorno would be hostile toward. As detailed in the poem, and as tantamount to the movement of factory machines, boredom can maintain a level of objectivity that Adorno would desire. Moreover, if boredom is, in fact, an adequate way of linking modern individuals and collectives, it could be a construct useful for unlocking what in ``every authentic artwork is internally revolutionary''.\footnote{Adorno, \textit{AT}, 228.}

\section{References}
%\addcontentsline{toc}{section}{Works Cited}
\begin{singlespacing}
\begin{hangparas}{.25in}{1}
Adorno, Theodor. \textit{Aesthetic Theory}. Translated and edited by Robert Hullot-Kentor. Minneapolis, MN: University of Minnesota Press, 1997.

---. ``The Essay as Form''. In \textit{Notes to Literature, Volume 1}. Edited by Ralph Tiedemann. Translated by Shierry Weber Nicholsen. New York: Columbia University Press, 1991.

---. ``On Subject and Object''. In \textit{Critical Models}. Translated by Henry Pickford. New York: Columbia University Press, 1999.

---. ``On the Use of Foreign Words''. In \textit{Notes to Literature, Volume 2}. Edited by Ralph Tiedemann. Translated by Shierry Weber Nicholsen. New York: Columbia University Press, 1992.

Baudelaire, Charles. ``\`A une passante''. In \textit{Fleurs du Mal (1861 Edition)}. Last modified 2018. \url{https://fleursdumal.org/poem/224}.

---. ``Au Lecteur''. In \textit{Fleurs du Mal (1861 Edition)}. Last modified 2018. \url{https://fleursdumal.org/poem/099}.

Benjamin, Walter. ``Some Motifs in Baudelaire''. In \textit{Illuminations: Essays and Reflections}. Translated by Harry Zohn. Edited by Hannah Arendt. New York: Schocken Books, 1968.

Hullot-Kentor, Robert. Introduction to \textit{Aesthetic Theory}. By Theodor Adorno. Translated and edited by Robert Hullot-Kentor. Minneapolis, MN: University of Minnesota Press, 1997.

Ngai, Sianne. \textit{Our Aesthetic Categories: Zany, Cute, Interesting}. Cambridge, MA: Harvard University Press, 2012.

Plato. \textit{Phaedrus}. Translated by Alexander Nehamas and Paul Woodruff. Indianapolis, IN: Hackett Publishing, 1995.
\end{hangparas}
\end{singlespacing}

\newpage
\appendix
\section{Baudelaire, ``Au Lecteur''}
\footnote{Charles Baudelaire, ``Au Lecteur'' in \textit{Fleurs du Mal (1861 Edition)}, last modified 2018, \url{https://fleursdumal.org/poem/099}.}

\begin{singlespacing}
\begin{verse}
La sottise, l'erreur, le p\'ech\'e, la l\'esine,\\
Occupent nos esprits et travaillent nos corps,\\
Et nous alimentons nos aimables remords,\\
Comme les mendiants nourrissent leur vermine.
\vspace{\baselineskip}

Nos p\'ech\'es sont t\^etus, nos repentirs sont l\^aches;\\
Nous nous faisons payer grassement nos aveux,\\
Et nous rentrons gaiement dans le chemin bourbeux,\\
Croyant par de vils pleurs laver toutes nos taches.
\vspace{\baselineskip}

Sur l'oreiller du mal c'est Satan Trism\'egiste\\
Qui berce longuement notre esprit enchant\'e,\\
Et le riche m\'etal de notre volont\'e\\
Est tout vaporis\'e par ce savant chimiste.
\vspace{\baselineskip}

C'est le Diable qui tient les fils qui nous remuent!\\
Aux objets r\'epugnants nous trouvons des appas;\\
Chaque jour vers l'Enfer nous descendons d'un pas,\\
Sans horreur, \`a travers des t\'en\`ebres qui puent.
\vspace{\baselineskip}

Ainsi qu'un d\'ebauch\'e pauvre qui baise et mange\\
Le sein martyris\'e d'une antique catin,\\
Nous volons au passage un plaisir clandestin\\
Que nous pressons bien fort comme une vieille orange.
\vspace{\baselineskip}

Serr\'e, fourmillant, comme un million d'helminthes,\\
Dans nos cerveaux ribote un peuple de D\'emons,\\
Et, quand nous respirons, la Mort dans nos poumons\\
Descend, fleuve invisible, avec de sourdes plaintes.
\vspace{\baselineskip}

Si le viol, le poison, le poignard, l'incendie,\\
N'ont pas encor brod\'e de leurs plaisants dessins\\
Le canevas banal de nos piteux destins,\\
C'est que notre \^ame, h\'elas! n'est pas assez hardie.
\vspace{\baselineskip}

Mais parmi les chacals, les panth\`eres, les lices,\\
Les singes, les scorpions, les vautours, les serpents,\\
Les monstres glapissants, hurlants, grognants, rampants,\\
Dans la m\'enagerie inf\^ame de nos vices,
\vspace{\baselineskip}

II en est un plus laid, plus m\'echant, plus immonde!\\
Quoiqu'il ne pousse ni grands gestes ni grands cris,\\
Il ferait volontiers de la terre un d\'ebris\\
Et dans un b\^aillement avalerait le monde;
\vspace{\baselineskip}

C'est l'Ennui! L'oeil charg\'e d'un pleur involontaire,\\
II r\^eve d'\'echafauds en fumant son houka.\\
Tu le connais, lecteur, ce monstre d\'elicat,\\
--- Hypocrite lecteur, --- mon semblable, --- mon fr\`ere!
\end{verse}
\end{singlespacing}

\end{document}
